\section{Preliminary Results}
\label{sec:results}

\noindent\textbf{Dataset and setup.}
We draw our image pool from SPECS (Street Perception Evaluation
Considering Socioeconomics)~\cite{quintana2025}, a demographically
balanced perception survey in which 1{,}000 participants from five
countries rated street-view scenes on ten perceptual indicators---six
traditional (safe, lively, wealthy, beautiful, boring, depressing) and
four new (live nearby, walk, cycle, green).
SPECS provides per-image perception scores grounded in diverse human
judgements, making it a principled source for studying which visual
levers shift a specific attribute.
We sample $N{=}50$ scenes from five SPECS cities (Amsterdam, Abuja,
San Francisco, Santiago, Singapore; 10 per city), stratified jointly by
baseline safety score and visual complexity so that the pilot spans the
full range of starting conditions---from low-safety scenes where
intervention is most needed to high-safety scenes that may offer fewer
actionable levers.
Each scene allows up to $K{=}5$ lever candidates with a stochastic
budget of $T{=}5$ generation attempts per candidate. In this paper, we
instantiate the framework with Qwen-Image-Edit as the prompt-only
image editing model and GPT-5.4 as the LLM-as-judge critic to
establish a reproducible lower-bound baseline.
Validity auditing serves as an eligibility gate throughout this section:
all effect summaries are computed only on edits retained in $V(x)$.

\noindent\textbf{Auxiliary directional shifts.}
Among the 177 valid edits, the auxiliary scorer produces a mean shift of
$+0.366$ (95\% CI $[+0.199, +0.537]$, median $+0.184$, range
$[-3.624, +3.842]$). Using threshold $\theta_{\mathrm{aux}}{=}0.1$,
95 of 177 valid edits fall in $E_{\mathrm{aux}}(x,a)$, corresponding to
a mean of 1.90 auxiliary-promising levers per scene
(95\% CI $[1.476, 2.324]$); 40 of the 50 scenes retain at least one
$\theta_{\mathrm{aux}}{=}0.1$ auxiliary-promising lever and 24 retain multiple.
In this pilot, the largest retained auxiliary shifts are observed for
lane-marking repainting, crosswalk repainting, and localized greenery
addition.
At the same time, some valid edits remain neutral or negative, reinforcing that
valid realization is necessary for analysis but not sufficient for
positive auxiliary movement.

\begin{table}[t]
  \centering
  \caption{Lever-type results grouped by intervention family.
  Mean~$\Delta_{\mathrm{aux}}$ and $\Delta_{\mathrm{aux}}$ [95\% CI]
  are computed over valid edits only; confidence intervals are
  bootstrap percentile intervals over the retained set.
  ``--'' denotes concepts not proposed in this scene pool.}
  \label{tab:levers}
  \scriptsize
  \resizebox{\columnwidth}{!}{%
  \begin{tabular}{p{0.43\linewidth}ccc}
    \toprule
    Family / Lever concept & Valid & Mean $\Delta_{\mathrm{aux}}$ & $\Delta_{\mathrm{aux}}$ [95\% CI] \\
    \midrule
    \textit{Physical Maintenance} & & & \\
    \quad Graffiti removal        & 4  & $+$0.396 & [$+$0.121, $+$0.670] \\
    \quad Litter removal          & 25 & $+$0.296 & [$-$0.128, $+$0.743] \\
    \quad Facade repair           & 4  & $+$0.519 & [$+$0.105, $+$0.949] \\
    \quad Surface cleaning        & 38 & $+$0.352 & [$+$0.026, $+$0.684] \\
    \quad\textbf{Family total}    & 71 & $+$0.344 & [$+$0.107, $+$0.584] \\
    \midrule
    \textit{Environmental Amenity} & & & \\
    \quad Localized greenery add.\  & 32 & $+$0.650 & [$+$0.241, $+$1.070] \\
    \quad Lighting repair           & 8  & $-$0.232 & [$-$0.598, $+$0.173] \\
    \quad Tree canopy management    & 14 & $-$0.245 & [$-$0.952, $+$0.446] \\
    \quad\textbf{Family total}      & 54 & $+$0.287 & [$-$0.041, $+$0.610] \\
    \midrule
    \textit{Visual Legibility} & & & \\
    \quad Signage decluttering         & 5 & $-$0.631 & [$-$1.294, $-$0.151] \\
    \quad Storefront transparency      & 2 & $+$0.958 & [$+$0.184, $+$1.732] \\
    \quad\textbf{Family total}         & 7 & $-$0.177 & [$-$0.927, $+$0.599] \\
    \midrule
    \textit{Mobility Infrastructure} & & & \\
    \quad Crosswalk repainting       & 15 & $+$0.767 & [$+$0.255, $+$1.261] \\
    \quad Lane marking repainting    & 30 & $+$0.485 & [$+$0.045, $+$0.934] \\
    \quad\textbf{Family total}       & 45 & $+$0.579 & [$+$0.244, $+$0.923] \\
    \midrule
    \textbf{Overall}                 & 177 & $+$0.366 & [$+$0.199, $+$0.537] \\
    \bottomrule
  \end{tabular}
  }
\end{table}

Table~\ref{tab:levers} shows that the main family-level signal lies in
the $\Delta_{\mathrm{aux}}$ column. Mobility Infrastructure yields the
largest mean auxiliary shift in this pilot ($+0.579$, 95\% CI $[+0.244, +0.923]$),
with crosswalk repainting ($+0.767$, 95\% CI $[+0.255, +1.261]$) and
lane-marking repainting ($+0.485$, 95\% CI $[+0.045, +0.934]$) leading
the study. Physical Maintenance is the broadest consistently positive
family ($+0.344$, 95\% CI $[+0.107, +0.584]$), driven by surface
cleaning, litter removal, facade repair, and graffiti removal.
Environmental Amenity is more heterogeneous: localized greenery
addition is strongly positive ($+0.650$, 95\% CI $[+0.241, +1.070]$),
whereas lighting repair and tree canopy management remain weak or
negative on average. Visual Legibility remains small-sample and mixed
in sign, so its directional pattern is not yet stable; the negative mean
for signage decluttering ($-0.631$, 95\% CI $[-1.294, -0.151]$) may
reflect the auxiliary model associating commercial signage density with
activity and perceived safety in Place Pulse training data. The Valid
column is included to show the retained sample supporting each estimate;
proposal and valid-rate diagnostics are reported in
Appendix Tables~\ref{tab:family_appendix} and~\ref{tab:city_appendix},
so the main table here remains focused on directional proxy movement.

\noindent\textbf{Thresholded auxiliary-promising lever counts.}
Figure~\ref{fig:delta_cutoff_counts} is the main per-scene summary
because it counts levers that are both valid and directionally
promising under the auxiliary proxy at increasingly strict auxiliary thresholds. At the study
threshold $\theta_{\mathrm{aux}}{=}0.1$, Mobility Infrastructure
averages 0.56 auxiliary-promising levers per scene with mean retained shift
$+1.270$ (95\% CI $[+0.958, +1.583]$), Environmental Amenity averages
0.54 with mean $+1.168$ (95\% CI $[+0.821, +1.516]$), and
Physical Maintenance averages 0.76 with mean $+1.044$ (95\% CI
$[+0.793, +1.296]$). Visual Legibility averages only 0.04 levers per
scene, so its interval is necessarily wide. By city, San Francisco and
Santiago dominate the positive tail at $\theta_{\mathrm{aux}}{=}0.1$
with averages of 2.5 and 2.6 auxiliary-promising levers per scene respectively,
while Abuja averages 1.1. As the cutoff rises to $1.0$, the family
ranking compresses but Environmental Amenity and Mobility
Infrastructure still retain the largest positive tails in this pilot, and
Amsterdam, San Francisco, and Santiago remain the cities with the
highest high-delta leverage per scene.

\begin{table}[t]
  \centering
  \caption{City-level summary over the full $N{=}50$ run.
  Mean $\Delta_{\mathrm{aux}}$ and $\Delta_{\mathrm{aux}}$ [95\% CI]
  are computed over valid edits only. Coverage diagnostics are reported
  in Appendix Table~\ref{tab:city_appendix}.}
  \label{tab:cities}
  \scriptsize
  \resizebox{\columnwidth}{!}{%
  \begin{tabular}{lccc}
    \toprule
    City & Valid & Mean $\Delta_{\mathrm{aux}}$ & $\Delta_{\mathrm{aux}}$ [95\% CI] \\
    \midrule
    Amsterdam     & 40 & $+$0.400 & [$+$0.073, $+$0.764] \\
    Abuja         & 32 & $-$0.136 & [$-$0.441, $+$0.141] \\
    San Francisco & 36 & $+$0.704 & [$+$0.285, $+$1.105] \\
    Santiago      & 38 & $+$0.742 & [$+$0.398, $+$1.089] \\
    Singapore     & 31 & $-$0.012 & [$-$0.383, $+$0.352] \\
    \midrule
    \textbf{Overall} & 177 & $+$0.366 & [$+$0.199, $+$0.537] \\
    \bottomrule
  \end{tabular}
  }
\end{table}

\begin{figure*}[t]
  \centering
  \includegraphics[width=0.9\textwidth]{figures/figure_5_delta_cutoff_counts.png}
  \caption{
    Average number of retained levers per scene satisfying
    $\Delta_{\mathrm{aux}} \geq \tau$, split by lever family (left) and
    city (right). The annotation reports the overall mean
    $\Delta_{\mathrm{aux}}$ with 95\% CI and the mean number of
    $\theta_{\mathrm{aux}}{=}0.1$ auxiliary-promising levers per scene with
    95\% CI.
  }
  \label{fig:delta_cutoff_counts}
\end{figure*}

Table~\ref{tab:cities} is best read as supporting context rather than a
primary preliminary signal. Directional auxiliary movement is more
context-dependent than validity gating: Santiago and San Francisco have
the largest positive mean shifts, whereas Abuja and Singapore are near
zero or slightly negative on average. The valid-rate differences remain
informative as a feasibility diagnostic, with the appendix table
showing Amsterdam and Santiago highest on retained valid counts and
Singapore lowest.

\noindent\textbf{Baseline score and gate-qualified editability.}
Appendix Figure~\ref{fig:appendix_baseline_vs_valid} plots each scene's baseline safety
score $f_a(x)$ against its valid lever count $|V(x)|$. At $N{=}50$, no
meaningful monotonic relationship is observed (Spearman
$\rho{=}0.10$, $p{=}0.511$). In other words, once generation reliability
is no longer the binding constraint, scenes with lower baseline safety
do not obviously admit more valid levers than higher-baseline scenes.
We treat this as a secondary diagnostic about editability rather than a
main perception result and revisit the null pattern in
\S\ref{sec:discussion}.

\noindent\textbf{Validity gate.}
Of 250 proposed candidates, 177 pass the validity audit (rate 0.708);
all 50 scenes retain at least one valid lever. Detailed coverage and
failure analysis are in Appendix~\ref{sec:validity_appendix}.

\noindent\textbf{Qualitative examples.}
Appendix Figure~\ref{fig:appendix_qualitative} shows three positive auxiliary-shift
examples from distinct scenes: lane-marking repainting in Amsterdam,
facade repair in Singapore, and surface cleaning in
Santiago. All three are spatially constrained, visually
legible, and plausible as single-lever urban changes.
